\documentclass[12pt]{article}
\usepackage[utf8]{inputenc}
\usepackage[T1]{fontenc}
\usepackage{hyperref}
\title{Ts\textquoteright\"un-Lao Language Resources for K-12 Curriculum Development}
\author{Compiled by Claude Sonnet 5 (claude-sonnet-5) for Andr\'as Kornai}
\date{}
\begin{document}
\maketitle

\section{Introduction}
Ts\textquoteright\"un-Lao (ISO 639-3 tsl, Glottocode tsun1242; also called L\`ao B\'oc) is a Kra-Dai (Tai) language spoken in Lai Ch\^au Province, Vietnam. Wikipedia gives about 14,900 members of the ethnic group (2009 census) and about 10,000 speakers (1993 figure). Its subgrouping is uncertain. Very little educational material appears to exist, and no school use was found in the search results.

\section{Primary Educational Resources}
No primers, textbooks or school programmes for Ts\textquoteright\"un-Lao were found in search results.

\section{Linguistic and Reference Materials}
\begin{itemize}
\item Wikipedia overview of the language (classification, location, speaker figures): \url{https://en.wikipedia.org/wiki/Ts%CA%BB%C3%BCn-Lao_language}
\end{itemize}
No dictionary, grammar or orthography was found.

\section{Religious and Community Resources}
\begin{itemize}
\item Joshua Project profile of the Tsun-Lao in Vietnam. It lists the language as unwritten and reports that Bible translation has started: \url{https://joshuaproject.net/people_groups/15599}
\item Global Recordings Network page for tsl (audio listing; its contents were not inspected): \url{https://globalrecordings.net/en/language/tsl}
\end{itemize}

\section{Digital and Audio Resources}
The only audio lead is the Global Recordings Network page above. No corpora or software were found.

\section{Curriculum Development Context}
The language appears to be unwritten and to have no documented educational materials. Any curriculum would first need community consultation, an orthography and basic documentation. Vietnamese is the likely language of schooling.
\end{document}